\section{Exact source-off return and section throughput}\label{sec:return}
For $\sigma\in\{+1,-1\}$, define normalized characteristic fields, their densities and stocks by
\begin{equation}\label{eq:chars}
 w_\sigma=\frac{p+\sigma r_c}{\sqrt2},\qquad
 \rho_\sigma=\frac{w_\sigma^2}{2},\qquad
 E_\sigma(t)=\int_0^{L_c}\rho_\sigma\dd s.
\end{equation}
Adding and subtracting \eqref{eq:C1} gives the inherited transport equation. Multiplication by $w_\sigma$ gives its directional account:
\begin{align}
 (\partial_t+\sigma c\partial_s)w_\sigma&=g_c/\sqrt2,\label{eq:transport}\\
 \partial_t\rho_\sigma+\sigma c\partial_s\rho_\sigma
 &=w_\sigma g_c/\sqrt2.\label{eq:directional}
\end{align}
The product calculation has the distributional meaning specified in Section~\ref{sec:workaccount}. Inversion of the orthogonal transform gives
\begin{equation}\label{eq:netgross}
 e=\rho_++\rho_-,\qquad
 j=c(\rho_+-\rho_-),\qquad
 j_{\rm gross}:=c(\rho_++\rho_-)=ce,\qquad |j|\le j_{\rm gross}.
\end{equation}
The signed directional fluxes are $+c\rho_+$ and $-c\rho_-$. Their nonnegative magnitudes add to the gross flux, which is generally not $|j|$.

\begin{theorem}[Conditional source-off recurrence]\label{thm:return}
On an entirely source-off interval, $E_+$ and $E_-$ are separately conserved. If $[t,t+T_{\rm loop}]$ is entirely source-off, with
\begin{equation}\label{eq:period}
 T_{\rm loop}=L_c/c,
\end{equation}
then $(p,r_c)(\cdot,t+T_{\rm loop})=(p,r_c)(\cdot,t)$. The equality holds in $L^2$ for mild preparation and in the continuous spatial representatives at $H^1$ regularity. $T_{\rm loop}$ is a period, not necessarily the smallest one.
\end{theorem}
\begin{proof}
The inherited characteristic representation \cite[well-posedness theorem]{pairedII} specializes, on a source-off interval starting at $a$, to
\begin{equation}\label{eq:translation}
 w_\sigma(s,t)=w_\sigma(s-\sigma c(t-a),a).
\end{equation}
Thus the $+$ profile travels toward increasing $s$ and the $-$ profile toward decreasing $s$, both at speed $c$. Periodic translation is an $L^2$ isometry, so it preserves each squared norm and hence each $E_\sigma$. Translation through $-\sigma cT_{\rm loop}=-\sigma L_c$ is the identity on the periodic domain. Inverting the characteristic transform proves recurrence of $(p,r_c)$. No newly evaluated wave solution is needed.
\end{proof}

\begin{theorem}[One-circuit section account]\label{thm:throughput}
At a declared fixed cut $s_*$, under the trace convention in Section~\ref{sec:workaccount}, define
\begin{equation}\label{eq:Q}
 Q_\sigma(s_*;a,b)=\int_a^b c\rho_\sigma(s_*,u)\dd u.
\end{equation}
If the entire interval $[t,t+T_{\rm loop}]$ is source-off, then
\begin{align}
 Q_\sigma(s_*;t,t+T_{\rm loop})&=E_\sigma(t),\label{eq:Qdirection}\\
 Q_{\rm net}:=\int_t^{t+T_{\rm loop}}j(s_*,u)\dd u
 &=E_+(t)-E_-(t),\label{eq:Qnet}\\
 Q_{\rm gross}:=\int_t^{t+T_{\rm loop}}j_{\rm gross}(s_*,u)\dd u
 &=E_+(t)+E_-(t)=E(t).\label{eq:Qgross}
\end{align}
\end{theorem}
\begin{proof}
Apply \eqref{eq:translation} with starting time $t$ in \eqref{eq:Q}. The change of variables $x=s_*-\sigma c(u-t)$ covers the circle once in time $L_c/c$. For $\sigma=+1$ the reversed limits cancel the minus sign from $\dd x=-c\dd u$; for $\sigma=-1$ the orientation is increasing. In either case periodicity gives
\[
 \int_t^{t+L_c/c}c\rho_\sigma(s_*-\sigma c(u-t),t)\dd u
 =\int_0^{L_c}\rho_\sigma(x,t)\dd x.
\]
This proves \eqref{eq:Qdirection}. The difference and sum, using \eqref{eq:netgross}, give \eqref{eq:Qnet} and \eqref{eq:Qgross}.
\end{proof}

The theorem distinguishes counterpropagation from net transport. Equal positive directional stocks have zero full-circuit net cut amount and positive gross throughput. It does not require a preferred direction or nonzero flux for every preparation.

Over $m$ whole source-off circuits, $m$ a positive integer, adding \eqref{eq:Qgross} gives gross throughput $mE$, while storage remains $E$. This recounts circulating storage, not $mE$ of new emission. A statement that all returns requires a declared profile or same-direction section-turnover meaning and a positive denominator. Neither theorem identifies tagged particles, a launch cohort in a superposed driven field, a native-state recurrence, or receipt at the spatial center.

\begin{corollary}[Zero source-off total-storage depletion]\label{cor:depletion}
If $E(t)>0$ and $[t,t+T_{\rm loop}]$ is source-off, then
\begin{equation}\label{eq:depletion}
 d_E=\frac{E(t)-E(t+T_{\rm loop})}{E(t)}=0.
\end{equation}
At $E(t)=0$ this ratio is undefined, while conservation still holds.
\end{corollary}
\begin{proof}
The sum $E=E_++E_-$ is conserved. A proposed law
$E(t+T_{\rm loop})=(1-d)E(t)$ with $d>0$ would therefore contradict conservation when $E(t)>0$.
\end{proof}
The same storage obstruction holds on every source-off interval, not only a full circuit. It is a theorem about unmodified C1, whose inspected equations contain no intrinsic sink, damping or leakage branch. It is not a physical measurement showing that the author's intended circulation is lossless. The source's statement that autonomous circulation has not been established concerns the missing physical and reciprocal core interpretation \cite[conclusion]{pairedII}; it does not negate homogeneous mathematical profile recurrence.
